\begin{exercises}
  \recommended \item 
    判断下列矩阵是否行等价。
    \begin{exparts*}
       \partsitem \(
           \begin{mat}[r]
             1  &2  \\
             4  &8
           \end{mat}, 
           \begin{mat}[r]
             0  &1  \\
             1  &2
           \end{mat} \)
       \partsitem \(
           \begin{mat}[r]
             1  &0  &2  \\
             3  &-1 &1  \\
             5  &-1 &5
           \end{mat},  
           \begin{mat}[r]
             1  &0  &2  \\
             0  &2  &10 \\
             2  &0  &4
           \end{mat} \)
       \partsitem \(
           \begin{mat}[r]
             2  &1  &-1 \\
             1  &1  &0  \\
             4  &3  &-1
           \end{mat},  
           \begin{mat}[r]
             1  &0  &2  \\
             0  &2  &10 \\
           \end{mat} \)
       \partsitem \(
           \begin{mat}[r]
             1  &1  &1  \\
            -1  &2  &2
           \end{mat},  
           \begin{mat}[r]
             0  &3  &-1 \\
             2  &2  &5
           \end{mat} \)
       \partsitem \(
           \begin{mat}[r]
             1  &1  &1  \\
             0  &0  &3
           \end{mat},  
           \begin{mat}[r]
             0  &1  &2  \\
             1  &-1 &1
           \end{mat} \)
    \end{exparts*}
    \begin{answer}
      把每个矩阵都化成简化阶梯形，然后进行比较。
      \begin{exparts}
        \partsitem 第一个给出
          \begin{equation*}
            \grstep{-4\rho_1+\rho_2}
            \begin{mat}[r]
              1  &2  \\
              0  &0
            \end{mat}
          \end{equation*}
          而第二个给出
          \begin{equation*}
            \grstep{\rho_1\leftrightarrow\rho_2}
            \begin{mat}[r]
              1  &2  \\
              0  &1
            \end{mat}
            \grstep{-2\rho_2+\rho_1}
            \begin{mat}[r]
              1  &0  \\
              0  &1
            \end{mat}
          \end{equation*}
          两个简化阶梯形矩阵并不相同，因此
          原来的矩阵不行等价。
        \partsitem 第一个如下。
          \begin{equation*}
            \grstep[-5\rho_1+\rho_3]{-3\rho_1+\rho_2}
            \begin{mat}[r]
              1  &0  &2  \\
              0  &-1 &-5 \\
              0  &-1 &-5
            \end{mat}
            \grstep{-\rho_2+\rho_3}
            \begin{mat}[r]
              1  &0  &2  \\
              0  &-1 &-5 \\
              0  &0  &0
            \end{mat}
            \grstep{-\rho_2}
            \begin{mat}[r]
              1  &0  &2  \\
              0  &1  &5  \\
              0  &0  &0
            \end{mat}
          \end{equation*}
          第二个如下。
          \begin{equation*}
            \grstep{-2\rho_1+\rho_3}
            \begin{mat}[r]
              1  &0  &2  \\
              0  &2  &10 \\
              0  &0  &0
            \end{mat}
            \grstep{(1/2)\rho_2}
            \begin{mat}[r]
              1  &0  &2  \\
              0  &1  &5  \\
              0  &0  &0
            \end{mat}
          \end{equation*}
          这两个矩阵行等价。
        \partsitem 这两个矩阵不行等价，因为它们的
          尺寸不同。
        \partsitem 第一个，
          \begin{equation*}
            \grstep{\rho_1+\rho_2}
            \begin{mat}[r]
              1  &1  &1  \\
              0  &3  &3
            \end{mat}
            \grstep{(1/3)\rho_2}
            \begin{mat}[r]
              1  &1  &1  \\
              0  &1  &1
            \end{mat}
            \grstep{-\rho_2+\rho_1}
            \begin{mat}[r]
              1  &0  &0  \\
              0  &1  &1
            \end{mat}
          \end{equation*}
          而第二个，
          \begin{equation*}
            \grstep{\rho_1\leftrightarrow\rho_2}
            \begin{mat}[r]
              2  &2  &5  \\
              0  &3  &-1
            \end{mat}
            \grstep[(1/3)\rho_2]{(1/2)\rho_1}
            \begin{mat}[r]
              1  &1  &5/2 \\
              0  &1  &-1/3
            \end{mat}
            \grstep{-\rho_2+\rho_1}
            \begin{mat}[r]
              1  &0  &17/6 \\
              0  &1  &-1/3
            \end{mat}
          \end{equation*}
          这两个矩阵不行等价。
        \partsitem 这里第一个是
          \begin{equation*}
            \grstep{(1/3)\rho_2}
            \begin{mat}[r]
              1  &1  &1  \\
              0  &0  &1
            \end{mat}
            \grstep{-\rho_2+\rho_1}
            \begin{mat}[r]
              1  &1  &0  \\
              0  &0  &1
            \end{mat}
          \end{equation*}
          而这是第二个。
          \begin{equation*}
            \grstep{\rho_1\leftrightarrow\rho_2}
            \begin{mat}[r]
              1  &-1 &1  \\
              0  &1  &2
            \end{mat}
            \grstep{\rho_2+\rho_1}
            \begin{mat}[r]
              1  &0  &3  \\
              0  &1  &2
            \end{mat}
          \end{equation*}
          这两个矩阵不行等价。
       \end{exparts}  
     \end{answer}
  \item 
    这些矩阵中哪些相互行等价？
    \begin{exparts*}
      \partsitem
         \(\begin{mat}
             1 &3 \\
             2 &4
           \end{mat}\)
      \partsitem
         \(\begin{mat}
             1 &5 \\
             2 &10
           \end{mat}\)
      \partsitem
         \(\begin{mat}
             1 &-1 \\
             3 &0
           \end{mat}\)
      \partsitem
         \(\begin{mat}
             2 &6 \\
             4 &10
           \end{mat}\)
      \partsitem
         \(\begin{mat}
             0  &1 \\
             -1 &0
           \end{mat}\)
      \partsitem
         \(\begin{mat}
             3 &3 \\
             2 &2
           \end{mat}\)
    \end{exparts*}
    \begin{answer}
      对每一个作高斯–若尔当消元。
      两个矩阵行等价当且仅当它们具有相同的
      简化阶梯形。
      下面是每个矩阵的简化形式。 
      % sage: M = matrix(QQ, [[1,3], [2,4]])
      % sage: gauss_jordan(M)
      % [1 3]
      % [2 4]
      %  take -2 times row 1 plus row 2
      % [ 1  3]
      % [ 0 -2]
      %  take -1/2 times row 2
      % [1 3]
      % [0 1]
      %  take -3 times row 2 plus row 1
      % [1 0]
      % [0 1]
      % sage: M = matrix(QQ, [[1,5], [2,10]])
      % sage: gauss_jordan(M)
      % [ 1  5]
      % [ 2 10]
      %  take -2 times row 1 plus row 2
      % [1 5]
      % [0 0]
      % sage: M = matrix(QQ, [[1,-1], [3,0]])
      % sage: gauss_jordan(M)
      % [ 1 -1]
      % [ 3  0]
      %  take -3 times row 1 plus row 2
      % [ 1 -1]
      % [ 0  3]
      %  take 1/3 times row 2
      % [ 1 -1]
      % [ 0  1]
      %  take 1 times row 2 plus row 1
      % [1 0]
      % [0 1]
      % sage: M = matrix(QQ, [[2,6], [4,10]])
      % sage: gauss_jordan(M)
      % [ 2  6]
      % [ 4 10]
      %  take -2 times row 1 plus row 2
      % [ 2  6]
      % [ 0 -2]
      %  take 1/2 times row 1
      %  take -1/2 times row 2
      % [1 3]
      % [0 1]
      %  take -3 times row 2 plus row 1
      % [1 0]
      % [0 1]
      % sage: M = matrix(QQ, [[0,1], [-1,0]])
      % sage: gauss_jordan(M)
      % [ 0  1]
      % [-1  0]
      %  swap row 1 with row 2
      % [-1  0]
      % [ 0  1]
      %  take -1 times row 1
      % [1 0]
      % [0 1]
      % sage: M = matrix(QQ, [[3,3], [2,2]])
      % sage: gauss_jordan(M)
      % [3 3]
      % [2 2]
      %  take -2/3 times row 1 plus row 2
      % [3 3]
      % [0 0]
      %  take 1/3 times row 1
      % [1 1]
      % [0 0]
      \begin{exparts*}
        \partsitem
          \(\begin{mat}
              1  &0  \\
              0  &1
            \end{mat} \)
        \partsitem
          \(\begin{mat}
              1  &5  \\
              0  &0
            \end{mat} \)
        \partsitem
          \(\begin{mat}
             1   &0  \\
             0   &1
            \end{mat} \)
        \partsitem
          \(\begin{mat}
             1   &0  \\
             0   &1
            \end{mat} \)
        \partsitem
          \(\begin{mat}
             1   &0  \\
             0   &1
            \end{mat} \)
        \partsitem
          \(\begin{mat}
             1   &1  \\
             0   &0
            \end{mat} \)
      \end{exparts*}
    \end{answer}
 \item 写出另外三个与所给矩阵行等价的矩阵。
   \begin{exparts*}
     \partsitem 
       \(\begin{mat}
           1 &3 \\
           4 &-1
         \end{mat}\)
     \partsitem
       \(\begin{mat}
           0 &1 &2 \\
           1 &1 &1  \\
           2 &3 &4
         \end{mat}\)
   \end{exparts*}
   \begin{answer}
     对每个小题，只需对起始矩阵作一些行
     变换即可。
      \begin{exparts}
        \partsitem
          把第一行乘以~$3$，得到下式。
          \begin{equation*}
            \begin{mat}
              3 &9  \\
              4 &-1
            \end{mat}
          \end{equation*}
          （这个行变换的选择并没有什么讲究；它 
          只是脑海中想到的第一个而已。）
          另外两个行变换是行对换 $\rho_1\leftrightarrow\rho_2$ 
          和倍加 $\rho_1+\rho_2$。
          \begin{equation*}
            \begin{mat}
              4 &-1  \\
              1 &3
            \end{mat}
            \quad
            \begin{mat}
              1 &3  \\
              5 &2
            \end{mat}
          \end{equation*}
        \partsitem 作同样的三个任意行变换，得到
          这三个。
          \begin{equation*}
           \begin{mat} 
            0 &3 &6 \\
            1 &1 &1  \\
            2 &3 &4  
           \end{mat}
           \quad         
           \begin{mat} 
            1 &1 &1  \\
            0 &1 &2 \\
            2 &3 &4  
           \end{mat}
           \quad         
           \begin{mat} 
            0 &1 &2 \\
            1 &2 &3  \\
            2 &3 &4  
           \end{mat}
           \quad         
          \end{equation*}
      \end{exparts}
    \end{answer}
  \recommended \item 对这个矩阵施行高斯消元法。
    把最终矩阵的每一行表示成起始矩阵
    各行的线性组合。
    \begin{equation*}
      \begin{mat}
        1 &2   &1 \\
        3 &-1  &0 \\
        0 &4   &0
      \end{mat}
    \end{equation*}
    \begin{answer}
      高斯化简是常规计算。
      % sage: load("gauss_method.sage")
      % sage: M = matrix (QQ, [[1,2,1], [3,-1,0], [0,4,0]])
      % sage: gauss_method(M)
      % [ 1  2  1]
      % [ 3 -1  0]
      % [ 0  4  0]
      %  take -3 times row 1 plus row 2
      % [ 1  2  1]
      % [ 0 -7 -3]
      % [ 0  4  0]
      %  take 4/7 times row 2 plus row 3
      % [    1     2     1]
      % [    0    -7    -3]
      % [    0     0 -12/7]
      \begin{equation*}
        \begin{mat}
          1 &2   &1 \\
          3 &-1  &0 \\
          0 &4   &0
        \end{mat}
        \grstep{-3\rho_1+\rho_2}    
        \begin{mat}
          1 &2   &1 \\
          0 &-7  &-3 \\
          0 &4   &0
        \end{mat}
        \grstep{(4/7)\rho_2+\rho_3}    
        \begin{mat}
          1 &2   &1 \\
          0 &-7  &-3 \\
          0 &0   &-12/7
        \end{mat}
      \end{equation*}
      把这三个矩阵分别记为 $A$、$D$ 和~$B$，则有下式。
      \begin{align*}
        \begin{mat}
          \alpha_1 \\
          \alpha_2 \\
          \alpha_3
        \end{mat}
        &\grstep{-3\rho_1+\rho_2}    
        \begin{mat}
          \delta_1=\alpha_1 \\
          \delta_2=-3\alpha_1+\alpha_2 \\
          \delta_3=\alpha_3
        \end{mat}                                   \\
        & \grstep{(4/7)\rho_2+\rho_3}
        \begin{mat}
          \beta_1=\alpha_1  \\
          \beta_2=-3\alpha_1+\alpha_2 \\
          \beta_3=-(12/7)\alpha_1+(4/7)\alpha_2+\alpha_3
        \end{mat}
      \end{align*}
    \end{answer}
  \item 
     描述
     \nearbyexample{ex:RowEqClassTwoTwoMats} 所展示的各个等价类中的矩阵。
     \begin{answer}
       首先，与全
       \( 0 \) 矩阵行等价的矩阵只有它自己（因为行变换对它没有效果）。

       其次，能化简为 
       \begin{equation*}
         \begin{mat}
           1  &a  \\
           0  &0
         \end{mat}
       \end{equation*}
       的矩阵形如
       \begin{equation*}
         \begin{mat}
           b  &ba \\
           c  &ca
         \end{mat}
       \end{equation*}
       （其中 \( a,b,c\in\Re \)，且 \(b\) 与 \(c\) 不全为零）。  

       再次，能化简为 
       \begin{equation*}
         \begin{mat}[r]
           0  &1  \\
           0  &0
         \end{mat}
       \end{equation*}
       的矩阵形如
       \begin{equation*}
         \begin{mat}
           0  &a \\
           0  &b
         \end{mat}
       \end{equation*}
       （其中 \( a,b\in\Re \)，且不全为零）。  

       最后，能化简为 
       \begin{equation*}
         \begin{mat}[r]
           1  &0  \\
           0  &1
         \end{mat}
       \end{equation*}
       的矩阵是非奇异的矩阵。
       这是因为以这个矩阵为系数矩阵的线性方程组将有唯一解，而这正是
       非奇异的定义。
       （换一种说法：它们不属于上面的任何一个等价类。）
     \end{answer}
  \item 
    描述与下列矩阵
    同一行等价类中的所有矩阵。
    \begin{exparts*}
       \partsitem  \(
           \begin{mat}[r]
             1  &0  \\
             0  &0
           \end{mat}  \)
       \partsitem  \(
           \begin{mat}[r]
             1  &2      \\
             2  &4
           \end{mat} \)
       \partsitem \(
           \begin{mat}[r]
             1  &1      \\
             1  &3
           \end{mat} \)
    \end{exparts*}
    \begin{answer}
      \begin{exparts}
        \partsitem 它们形如
          \begin{equation*}
            \begin{mat}
              a  &0  \\
              b  &0
            \end{mat}
          \end{equation*}
          其中 \( a,b\in\Re \) 至少有一个不为零。
        \partsitem 它们具有如下形式 
          \begin{equation*}
            \begin{mat}
             a  &2a \\
             b  &2b
            \end{mat}
          \end{equation*}
          其中 \( a,b\in\Re \) 至少有一个不为零。
        \partsitem 所给矩阵是非奇异的。
          所以该行等价类由所有非奇异的 $\nbyn{2}$
          矩阵组成。
          （给出公式的话，它们形如
          \begin{equation*}
            \begin{mat}
              a  &b  \\
              c  &d
            \end{mat}
          \end{equation*}
          其中 \( a,b,c,d\in\Re \)）且 \( ad-bc\neq 0 \)。
          我们将在第四章看到，这个公式 
          判定了 \( \nbyn{2} \) 矩阵
          何时非奇异。）
      \end{exparts}  
    \end{answer}
  \item 
    行等价类有多少个？
    \begin{answer}
       有无穷多个。
       例如，在 
       \begin{equation*}
         \begin{mat}
           1  &k  \\
           0  &0
         \end{mat}
       \end{equation*}
       中，每个 $k\in\Re$ 给出一个不同的等价类。  
    \end{answer}
  \item 
    行等价类能包含不同尺寸的矩阵吗？
    \begin{answer}
      不能。
      行变换不改变矩阵的尺寸。  
    \end{answer}
  \item 
    行等价类有多大？
    \begin{exparts} 
      \partsitem 证明：对任意全零矩阵，其等价类是有限的。
      \partsitem 还有别的等价类只含有限多个成员吗？
    \end{exparts}
    \begin{answer}
     \begin{exparts}
      \partsitem 对全零矩阵作行变换没有效果。
        因此每个这样的矩阵在它的行等价类中只有自己。  
      \partsitem 没有。
        任何非零元都可以倍乘。
     \end{exparts}
    \end{answer}
  \recommended \item 
    给出两个首主元位于相同列的
    简化阶梯形矩阵，
    但它们并不行等价。
    \begin{answer}
      这里有两个。
      \begin{equation*}
        \begin{mat}[r]
          1  &1  &0  \\
          0  &0  &1
        \end{mat}
        \quad\text{和}\quad
        \begin{mat}[r]
          1  &0  &0  \\
          0  &0  &1
        \end{mat}
      \end{equation*}  
     \end{answer}
  \recommended \item 
    证明任意两个 \( \nbyn{n} \) 非奇异矩阵
    是行等价的。
    任意两个奇异矩阵也行等价吗？
    \begin{answer}
      任意两个 \( \nbyn{n} \) 非奇异矩阵具有
      相同的简化阶梯形，即除对角线上
      为 \( 1 \) 之外全为 \( 0 \) 的那个矩阵。
      \begin{equation*}
        \begin{mat}
          1  &0  &       &0  \\
          0  &1  &       &0  \\
             &   &\ddots &   \\
          0  &0  &       &1
        \end{mat}
      \end{equation*}

      两个同尺寸的奇异矩阵未必行等价。
      例如，这两个 \( \nbyn{2} \) 奇异矩阵
      不行等价。
      \begin{equation*}
        \begin{mat}[r]
          1  &1  \\
          0  &0
        \end{mat}
        \quad\text{和}\quad
        \begin{mat}[r]
          1  &0  \\
          0  &0
        \end{mat}
      \end{equation*}  
    \end{answer}
  \recommended \item 
    描述包含下列矩阵的所有行等价类。
    \begin{exparts*}
      \partsitem \( \nbyn{2} \)~矩阵
      \partsitem \( \nbym{2}{3} \)~矩阵
      \partsitem \( \nbym{3}{2} \)~矩阵
      \partsitem \( \nbyn{3} \)~矩阵
    \end{exparts*}
    \begin{answer}
      由于每个等价类中有且仅有一个简化阶梯形矩阵，
      我们只需列出所有可能的简化阶梯形矩阵。

      这份清单参见 \nearbyexercise{exer:PossRedEchFrms} 的答案。 
    \end{answer}
  \item  
     \begin{exparts}
          \partsitem 证明：向量 $\vec{\beta}_0$ 是集合 $\set{\vec{\beta}_1,\ldots,\vec{\beta}_n}$
            中成员的线性组合当且仅当存在线性关系 
            $\zero=c_0\vec{\beta}_0+\cdots+c_n\vec{\beta}_n$
            其中 $c_0$ 不为零。
            （\textit{提示。}   当心 $\vec{\beta}_0=\zero$ 的情形。）
          \partsitem 利用它来简化 
             \nearbylemma{le:EchFormNoLinCombo} 的证明。   
       \end{exparts}
       \begin{answer}
          \begin{exparts}
           \partsitem 如果存在 $c_0$ 不为零的线性关系
             那么我们可以从两边减去 $c_0\vec{\beta}_0$ 再除以 $-c_0$，
             把 $\vec{\beta}_0$ 表示成其余向量的线性
             组合。
             （注：  
             如果集合中没有其他向量\Dash 即关系式
             是 $\zero=3\cdot\zero$\Dash 那么该命题仍然成立，因为
             按定义零向量是空向量集 
             的和。）

             反之，如果 $\vec{\beta}_0$ 是其余向量的组合 
             $\vec{\beta}_0=c_1\vec{\beta}_1+\dots+c_n\vec{\beta}_n$
             那么从两边减去
             $\vec{\beta}_0$ 就给出一个系数
             不全为零的线性关系；即
             $\vec{\beta}_0$ 前面的 $-1$。
           \partsitem 第一行不是其余各行的线性组合，
             理由
             如证明中所述：~在包含第一行首主元的那一列
             的分量等式中，唯一非零的元素
             是来自第一行的首主元，因此
             它的系数必为零。
             于是，由本题前面的小题，第一行与其
             余各行没有线性关系。

             因此，在考虑第二行能否与其余各行有线性 
             关系时，
             我们可以不考虑第一行。
             但这时刚才对第一行的论证也将适用于
             第二行。
             （也就是说，我们这里是用归纳法论证的。）             
         \end{exparts}
      \end{answer}
  % \item 
  %   Why, in the proof of \nearbytheorem{th:ReducedEchelonFormIsUnique},
  %   do we bother to restrict to the nonzero rows?
  %   Why not just stick to the relationship that we began with,
  %   $\beta_i=c_{i,1}\delta_1+\dots+c_{i,m}\delta_m$, with $m$ instead of $r$,
  %   and argue using it that the only nonzero coefficient
  %   is \( c_{i,i}  \), which is \( 1 \)?
  %   \begin{answer}
  %      The zero rows could have nonzero coefficients, and
  %      so the statement would not be true.
  %   \end{answer}
  \recommended \item 
   \cite{Trono}
   三位卡车司机走进一家路边咖啡馆。
   一位卡车司机买了四个三明治、一杯咖啡和十个 
   甜甜圈，花了 \$$8.45$。
   另一位司机买了三个三明治、一杯咖啡和七个
   甜甜圈，花了 \$$6.30$。
   第三位卡车司机买一个三明治、一杯咖啡和 
   一个甜甜圈要付多少钱？
    \begin{answer}
      我们知道 $4s+c+10d=8.45$ 与 $3s+c+7d=6.30$，而且想知道
      $s+c+d$ 是多少。
      幸运的是，$s+c+d$ 是 $4s+c+10d$ 与 $3s+c+7d$ 的线性组合。
      把未知的价格记为 $p$，我们有如下化简过程。
      \begin{align*}
        \begin{amat}{3}
          4  &1  &10  &8.45 \\
          3  &1  &7   &6.30 \\
          1  &1  &1   &p
        \end{amat}
        &\grstep[-(1/4)\rho_1+\rho_3]{-(3/4)\rho_1+\rho_2}
        \begin{amat}{3}
          4  &1    &10     &8.45      \\
          0  &1/4  &-1/2   &-0.037\,5 \\
          0  &3/4  &-3/2   &p-2.112\,5
        \end{amat}                                      \\
        &\grstep{-3\rho_2+\rho_3}
        \begin{amat}{3}
          4  &1    &10     &8.45      \\
          0  &1/4  &-1/2   &-0.037\,5 \\
          0  &0    &0      &p-2.00
        \end{amat}
     \end{align*}
     所付的价钱是 \$$2.00$。
   \end{answer}
   % \item
   %  The fact that Gaussian reduction disallows multiplication of
   %  a row by zero is needed for the proof of uniqueness of reduced echelon form,
   %  or else every matrix would
   %  be row equivalent to a matrix of all zeros.
   %  Where is it used?
   %  \begin{answer}
   %    If multiplication of a row by zero were allowed then
   %    \nearbylemma{le:EquivMatsSameForm}
   %    would not hold.
   %    That is, where
   %    \begin{equation*}
   %      \begin{mat}[r]
   %        1  &3  \\
   %        2  &1
   %      \end{mat}
   %      \grstep{0\rho_2}
   %      \begin{mat}[r]
   %        1  &3  \\
   %        0  &0
   %      \end{mat}
   %    \end{equation*}
   %    all the rows of the second matrix can be expressed as linear combinations
   %    of the rows of the first, but the converse does not hold.
   %    The second row of the first matrix is not a linear combination of the
   %    rows of the second matrix.  
   %  \end{answer}
  \item 
   线性组合引理说的是，对给定的线性方程组作
   高斯消元能得到哪些方程。
   \begin{enumerate}
     \item 写出一个不被这个方程组蕴含的方程。
       \begin{equation*}
         \begin{linsys}{2}
           3x  &+  &4y  &=  &8 \\
           2x  &+  & y  &=  &3 
         \end{linsys}
       \end{equation*}
     \item 从无解的方程组能导出任何方程吗？
   \end{enumerate}
   \begin{answer}
     \begin{enumerate}
        \item 一个容易的答案是这个。
          \begin{equation*}
            0=3
          \end{equation*}
          想要一个不那么抬杠的答案，求解该方程组：
          \begin{equation*}
            \begin{amat}[r]{2}
              3  &4  &8  \\
              2  &1   &3
            \end{amat}
            \grstep{-(2/3)\rho_1+\rho_2}
            \begin{amat}[r]{2}
              3  &4    &8    \\
              0  &-5/3 &-7/3
            \end{amat}
          \end{equation*}
          得到 \( y=7/5 \) 与 \( x=4/5 \)。
          于是任何不被 \( (7/5,4/5) \) 满足的方程都可以，
          例如 \( 5x+5y=10 \)。
        \item 从无解的方程组可以导出每一个方程。
          例如，下面演示如何从
          \( 0=5 \) 导出 \( 3x+2y=4 \)。
          首先，
          \begin{equation*}
            0=5
            \grstep{(3/5)\rho_1}
            0=3
            \grstep{x\rho_1}
            0=3x
          \end{equation*}
          （\( x=0 \) 情形的有效性需要另外论证，但很显然）。
          类似地，\( 0=2y \)。
          对 \( 0=4 \) 亦然。
          而现在，\( 0+0=0 \) 便给出 \( 3x+2y=4 \)。
     \end{enumerate}  
    \end{answer}
  \item 
    \cite{HoffmanKunze}
    把行等价的定义推广到线性方程组。
    在你的定义下，等价的方程组有相同的解集吗？
    \begin{answer}
      定义：若它们的增广
      矩阵行等价，则称线性方程组等价。
      等价方程组有相同解集的证明很容易。  
    \end{answer}
  \item 
    在这个矩阵
    \begin{equation*}
      \begin{mat}[r]
        1  &2  &3  \\
        3  &0  &3  \\
        1  &4  &5
      \end{mat}
    \end{equation*}
    中，第一列与第二列之和等于第三列。
    \begin{exparts}
      \partsitem 证明在任何行变换下这一点仍然成立。
      \partsitem 提出一个猜想。
      \partsitem 证明它成立。
    \end{exparts}
    \begin{answer}
      \begin{exparts}
        \partsitem 三种可能的行对换很容易， 
          三种可能的倍乘也是如此。
          六种可能的倍加之一是 \( k\rho_1+\rho_2 \)：
          \begin{equation*}
            \begin{mat}
              1           &2           &3  \\
              k\cdot 1+3  &k\cdot 2+0  &k\cdot 3+3  \\
              1           &4           &5
            \end{mat}
          \end{equation*}
          此时第一列与第二列之和仍等于第三列。
          其余五种倍加类似。
        \partsitem 显然的猜想是：行变换不改变
          列之间的线性关系。
        \partsitem 逐情形的
          证明按照第一小题给出的思路进行。
      \end{exparts}  
   \end{answer}
\end{exercises}
